\chapter{Gastrostomy and Jejunostomy Tube Placement}

\section*{Introduction \& Clinical Indications}
Gastrostomy and jejunostomy tubes provide enteral nutrition when oral intake is impossible or insufficient. A gastrostomy delivers feed to the stomach. A jejunal tube or gastrostomy with jejunal extension is considered when gastric feeding is not tolerated, gastric emptying is impaired, aspiration risk is high, or anatomy makes the gastric route unsuitable. Temporary nasogastric or nasojejunal feeding is generally preferred when enteral support is expected to last less than about four weeks.

Common indications include persistent dysphagia after neurologic disease, obstructing head and neck or esophageal disease, prolonged critical illness, and chronic disease with inadequate oral intake. Tube placement does not by itself prevent aspiration or improve outcomes in every condition; goals, prognosis, cognition, and alternatives should be discussed with the patient or surrogate.

\section*{Relevant Surgical Anatomy}
Safe gastrostomy requires an accessible anterior gastric wall apposed to the abdominal wall without intervening bowel, liver, or other organ. The stomach is supplied by branches of the celiac axis. Jejunal access follows the mesenteric vessels and requires attention to the small-bowel mesentery, abdominal wall, and any altered anatomy from prior surgery. The exact site is selected by endoscopic, radiologic, or surgical assessment rather than by a fixed surface measurement.

\section*{Preoperative Workup \& Preparation}
The team confirms the indication, expected duration of feeding, swallowing and aspiration assessment, prognosis, goals of care, and whether gastric or jejunal delivery is appropriate. Medication review includes anticoagulants and antiplatelet agents; peri-procedural management is individualized because percutaneous tube placement has meaningful bleeding risk. Examination and selective laboratory tests assess anemia, renal function, electrolytes, and coagulation risk. Imaging or contrast studies may be needed for altered anatomy, distension, ascites, or concern for interposed organs.

Informed consent covers bleeding, infection, visceral injury, leakage, peritonitis, tube malfunction, aspiration, need for revision, and procedure-related mortality. A single intravenous prophylactic antibiotic dose is used for PEG or PEJ according to local protocol. Uncorrectable coagulopathy, inability to safely oppose stomach and abdominal wall, active peritonitis, hemodynamic or respiratory instability, and an inconsistent treatment goal are reasons to defer or reconsider placement. Ascites and ventriculoperitoneal shunts increase risk but are not automatic absolute contraindications.

\section*{Surgical Technique \& Procedure}
\begin{enumerate}
  \item \textbf{Positioning and planning}: The patient is placed supine. For PEG, endoscopy confirms a safe puncture site using transillumination, indentation, and a safe-track technique.
  \item \textbf{Access}: After skin preparation and local anesthesia, the stomach is punctured under direct endoscopic, radiologic, or surgical guidance. A guidewire is placed and the tract is dilated as required.
  \item \textbf{Tube placement}: A PEG is commonly placed using the pull technique; a direct introducer technique is used when the pull route is unsuitable, such as severe esophageal obstruction or selected head and neck cancers. The internal and external retainers are positioned without excessive tension.
  \item \textbf{Jejunal access}: A jejunostomy or jejunal extension is positioned under the chosen endoscopic, radiologic, or surgical guidance, with confirmation of patency and location.
  \item \textbf{Completion}: The tube is secured, the site is dressed, and the tube's flush and feeding plan are documented.
\end{enumerate}

\section*{Complications \& Risks}
Early complications include peristomal infection, bleeding, visceral or bowel injury, leakage, aspiration, ileus, and peritonitis. Later complications include blockage, accidental dislodgement, buried bumper syndrome, granulation tissue, leakage, stomal narrowing, and recurrent aspiration. A displaced tube must not be blindly replaced through an immature tract; urgent clinical assessment is needed.

\section*{Postoperative Care \& Recovery}
The insertion site, abdominal examination, pain, vital signs, and tube position are checked. Feeding is started according to the procedural team's protocol once immediate complications have been excluded, and the dietitian sets the formula, rate, water flushes, and monitoring plan. Caregivers are taught hand hygiene, site care, tube flushing, medication administration, and what to do for blockage or dislodgement. Fever, worsening abdominal pain, peritoneal signs, persistent vomiting, bleeding, or leakage requires prompt assessment. The external bumper should remain appropriately loose; routine tube mobilization follows local protocol and device instructions.

\section*{Outcomes \& Prognosis}
Technical success is high in suitable candidates, but clinical benefit depends on the underlying disease, nutritional goals, prognosis, and complications. Placement does not guarantee improved survival, swallowing, or prevention of aspiration. Ongoing nutritional, functional, and goals-of-care review is part of safe long-term management.

\section*{Additional Clinical Considerations}
The route and duration of enteral support should be decided before the procedure whenever possible. A short-term nasogastric or nasojejunal tube may be sufficient during a reversible illness, whereas a percutaneous device is considered when feeding is expected to continue beyond roughly four to six weeks. The estimate is not a rigid rule: a patient may need earlier durable access when nasal tubes are repeatedly removed or poorly tolerated, and a percutaneous procedure may be inappropriate when recovery is expected to be brief or the goals of care do not include artificial nutrition.

Gastric feeding is usually preferred when the stomach empties adequately and the patient can protect the airway. Jejunal feeding may be useful in severe gastroparesis, gastric outlet obstruction, recurrent vomiting, or aspiration of gastric contents, but it does not eliminate aspiration because refluxed formula and oropharyngeal secretions can still be aspirated. Head-of-bed elevation when clinically safe, oral care, swallowing assessment, medication review, and treatment of reflux or delayed emptying remain important. A tube should not be presented as a substitute for dysphagia rehabilitation or as a guarantee against pneumonia.

Tube placement requires a safety pause. The team confirms the patient, indication, planned device, route, allergies, anticoagulant plan, relevant imaging, and rescue plan if access is difficult. Prior abdominal surgery, interposed colon or liver, severe ascites, portal hypertension, peritoneal carcinomatosis, and altered anatomy may favor radiologic or operative placement rather than blind or routine endoscopic access. A patient who cannot tolerate sedation or who has unstable respiratory or hemodynamic status may need stabilization or a different setting before placement.

For a newly placed nasoenteric tube, feeding and medication administration must wait until position is verified according to institutional policy. Radiographic confirmation is required for a blindly placed tube in many adult protocols; aspirate appearance and pH can provide supporting information but should not replace the required verification. Auscultation of an injected air bolus is not a reliable method for excluding pulmonary placement. The external length or marking is documented, and a change in that measurement, new coughing, respiratory distress, unexplained desaturation, or feeding intolerance prompts reassessment.

For a gastrostomy, the tract matures over time. Early accidental removal can cause leakage into the peritoneum, and blind replacement of a tube through an immature tract can be dangerous. If the tube comes out, the patient or caregiver should cover the site and contact the treating service rather than force another tube into the opening. A mature tract may still narrow quickly, so timely assessment is important. The team should document the tube type, French size, balloon or bumper details, external length, date of insertion, and instructions for replacement.

Feeding is advanced cautiously. The dietitian determines formula, caloric and protein targets, free-water needs, infusion method, and monitoring for refeeding syndrome in patients with prolonged inadequate intake. Electrolytes, glucose, hydration, weight, stool output, abdominal symptoms, and tolerance are followed. Medications should be checked for liquid formulation, compatibility, and interactions; sustained-release or enteric-coated products should not be crushed, and the tube should be flushed before and after medications according to local policy. The feeding port should not be used for substances that are incompatible with the prescribed route.

Long-term complications are often managed without removing the tube. Granulation tissue may need local treatment, leakage may reflect excessive tract tension or gastric pressure, and blockage may respond to protocolized flushing rather than forceful instrumentation. Persistent pain, buried bumper features, fever, peritoneal signs, bleeding, or new leakage around the tube requires clinician review. When oral intake becomes safe and adequate, removal is planned with the nutrition, swallowing, and procedural teams; the indication for placement, current goals, and risk of recurrent malnutrition are reassessed before removal.

\section*{Additional Follow-Up Considerations}
The patient and caregiver should receive a written tube record containing the device type, size, route, external length, insertion date, feeding plan, flush instructions, medication precautions, and emergency contact. A tube that is cracked, blocked, leaking, painful, or unexpectedly displaced should not be used until its position and integrity are assessed. Forceful flushing, blind replacement, and use of a tube for an unapproved medication or solution can cause injury.

Long-term review asks whether the tube is still meeting a meaningful nutritional or decompression goal. Weight, hydration, laboratory values, wound condition, swallowing function, oral intake, quality of life, and caregiver burden are considered. If the patient can eat safely and meet requirements, removal may be appropriate; if the indication persists, the tube and the care plan are reviewed rather than left unattended. Decisions should respect capacity, prior wishes, surrogate standards, and the overall goals of treatment.

\section*{Ethics and Goals of Care}
Artificial nutrition is a medical treatment, not an obligation. Capacity, prior wishes, prognosis, burdens, and the goals of the patient or legally authorized surrogate are reviewed. A tube may be appropriate for reversible malnutrition, rehabilitation, or long-term support in a patient who benefits from it, but it may not improve outcomes in every advanced illness. Decisions can be revisited as swallowing, function, and goals change.

\section*{Documentation and Safety}
The procedural record should identify the device, route, tip position, external length, fixation, feeding plan, and method used to confirm placement. A standardized handover prevents wrong-route feeding and medication errors. The patient or caregiver should demonstrate tube care before discharge and know that respiratory symptoms, new abdominal pain, leakage, bleeding, fever, or tube displacement require prompt clinical advice.

\section*{Practical Follow-Up}
The team monitors weight, hydration, electrolytes, feeding tolerance, tube position, skin condition, and the ongoing indication. Caregivers should not force a blocked tube or blindly replace a displaced tube. Any new abdominal pain, respiratory symptom, bleeding, fever, leakage, or unexpected external-length change requires clinical advice.

\section*{Safety Summary}
Enteral access is selected for a documented nutritional or decompression goal. Correct device identification and position verification precede feeding or medication administration, and the indication is revisited as the patient's condition changes.

\section*{Final Safety Note}
Tube type and position must be confirmed before use. A displaced, blocked, leaking, or painful tube should not be forcefully manipulated; the patient or caregiver should contact the responsible clinical service.

\section*{Completion Checklist}
Record the tube type, size, route, tip position, external length, feeding order, flush plan, and verification method. Do not feed through a displaced or uncertain tube; obtain clinical confirmation first.

\section*{Handoff}
The tube handoff should identify route, position verification, feeding and medication orders, external markings, and the emergency plan for displacement.

\section*{References}
\begin{enumerate}
  \item American Society for Gastrointestinal Endoscopy. Guideline on Gastrostomy Feeding Tubes. Gastrointestinal Endoscopy. 2025;101:25--35.
  \item European Society of Gastrointestinal Endoscopy. Endoscopic Management of Enteral Tubes in Adult Patients: Definitions and Indications. Endoscopy. 2021;53.
  \item European Society of Gastrointestinal Endoscopy. Endoscopic Management of Enteral Tubes in Adult Patients: Peri- and Post-procedural Management. Endoscopy. 2021.
\end{enumerate}

\clearpage
